\section{Conventions for the bilateral criterion}
\label{app:geometric-conventions}

We justify the precise version of Theorem~\ref{thm:bwgl}, including
positive image mass and bounded supports. The argument first transfers
our open-ball packing estimate to the metric coefficient used by
Bate, Hyde and Schul, and then turns one of their bi-Lipschitz pieces
into a Lipschitz image of the full parameter ball.

\begin{proof}[Proof of Theorem~\ref{thm:bwgl}]
The formulation of Bate, Hyde and Schul uses closed balls and maps
defined on all of $\R^n$ \cite[Theorem~1.0.4]{BateHydeSchul2023}.
If $\overline b_E$ denotes its closed-ball bilateral coefficient, then
\[
 \overline b_E(x,t)\leq2b_E(x,2t).
\]
Indeed, both closed-ball suprema at radius $t$ are bounded by the
open-ball suprema at radius $2t$ for the same plane. For
$R<\diam(E)/2$, enlarging the spatial ball to $B(a,2R)$ and substituting
$s=2t$ therefore transfers the Carleson hypothesis of
Theorem~\ref{thm:bwgl} to closed balls. If $D=\diam(E)<\infty$ and
$D/2\leq R<D$, cover $E$ by a
dimensionally bounded number of support-centered open balls of radius
$D/8$. On each enclosing radius-$D/4$ ball, the open-ball estimate
controls the transformed integral for $t<D/8$. For the remaining
$D/8\leq t<R$, the integral is at most
$(\log8)\H^n(E)\leq C D^n\leq C2^nR^n$.
Thus the closed-ball criterion applies with controlled constants in
both diameter regimes. AD regularity has the same open/closed-ball
equivalence, by doubling the radius and using the total-mass bound
when that doubled radius exceeds a finite diameter.

We next pass from Euclidean bilateral flatness to the metric coefficient
used in that source's corona decomposition. Its coefficient $\xi_E(a,t)$
\cite[Definition~3.1.3]{BateHydeSchul2023} minimizes, over maps from
$E\cap\overline B(a,t)$ to closed radius-$t$ balls in $n$-dimensional normed
spaces, the sum of two quantities: the largest pairwise distance
distortion divided by $t$, and the largest distance from a target-ball
point to the image divided by $t$. We claim
\[
 \xi_E(a,t)\leq36\overline b_E(a,2t).
\]
Choose an almost-fitting affine plane $L$ whose summed directed error
at radius $2t$ is at most $2\lambda t$, with $0<\lambda<1/12$.
Orthogonal projection, followed by translation, gives
$f(x)=\pi_L(x)-\pi_L(a)$ in the Euclidean closed radius-$t$ ball in the
direction of $L$. Its distance distortion is at most $4\lambda t$.
For a target vector $v$ with $|v|\leq t$, the point
$y=\pi_L(a)+(1-6\lambda)v$ satisfies
$|y-a|\leq(1-4\lambda)t$. The bilateral estimate supplies $z\in E$
with $|z-y|\leq2\lambda t$, so $z\in\overline B(a,t)$ and
$|f(z)-v|\leq8\lambda t$. Thus $\xi_E(a,t)\leq12\lambda$.
Letting $\lambda$ decrease to the infimum proves the claim when that
infimum is below $1/12$. Otherwise the constant map gives
$\xi_E(a,t)\leq3\leq36\overline b_E(a,2t)$.
The same doubling of scales and finite-diameter coarse-scale bound
proved above now transfer the Carleson estimate to $\xi_E$.
By \cite[Theorem~B, implication $(3)\Rightarrow(5)$]{BateHydeSchul2023},
$E$ therefore has the corona decomposition by normed spaces required
in that source's positive-mass construction, with controlled constants.

This corona decomposition supplies the bi-Lipschitz pieces in
\cite[Proposition~9.0.2 and Lemma~9.0.4]{BateHydeSchul2023}.
Choosing the omitted mass below half the lower AD bound gives a piece
$F\subset E\cap\overline B(a,s)$ with
$\H^n(F)\geq c s^n>0$ and a uniformly $K$-bi-Lipschitz map
$\phi:F\to\R^n$.
To see directly that a full parameter-ball image suffices, fix $p\in F$
and put $V=(\phi(F)-\phi(p))/(4K)$. Then
$V\subset B^n(0,s)$, and
$h(u)=\phi^{-1}(4Ku+\phi(p))$ on $V$ is $4K^2$-Lipschitz.
Coordinatewise Lipschitz extension gives a map on all of $\R^n$ with
constant at most $4\sqrt d K^2$ and image containing $F$.
Choose the preceding bi-Lipschitz piece at $s=r/2$ and restrict this map to the open ball
$B^n(0,r)$. Its image contains $F$, while
$\overline B(a,r/2)\subset B(a,r)$; the mass bound loses only the
factor $2^{-n}$. This proves precisely the open-ball formulation of
Theorem~\ref{thm:bwgl}.
\end{proof}
